\subsection{Group Averaging and Pauli Twirlin}
We illustrate some simple examples of Haar measure computations.
In particular, we are interested in first-order computations on discrete and continuous groups.
We approach this from the perspective of group averaging and twirling.
\begin{definition}[Group average]\label{def:group_average}
    Let $G$ be a group of order $N$ and $R$ a representation of the group.
    Define the {group average} to be
    \begin{equation}
        \langle X \rangle_{G} = \frac{1}{N} \sum_{g} R(g) X R(g)^{\dagger}.
    \end{equation}
\end{definition}

\begin{remark}[Assumptions]
    Consider a $n$-qubit system. Let $d = 2^n$, and suppose we seek to compute moments of order $k$.
    We assume for now that $d > k$ for the examples below.
\end{remark}

\subsection{Pauli twirling}
The end goal of this section is to prove the Irrep Group Averaging Corollary
\begin{equation}
    \langle X \rangle_{G} = \frac{1}{d} \text{Tr}[X] I,
\end{equation}
where $n_a = d$ is the dimension of the vector space of the representation.
To illustrate this, we first consider the group average of the single-qubit Pauli group on some arbitrary initial state $\rho$.
\begin{remark}[Pauli groups]\label{rm:pauli_group}
    We denote the single-qubit Pauli group as
    \begin{equation}
        G = \{\pm(i) \sigma_{x}, \pm(i) \sigma_y, \pm (i) \sigma_z, \pm (i) I\},
    \end{equation}
    where the Pauli matrices are given by
    \begin{equation}
        \sigma_0 = I = \begin{bmatrix} 1 & 0 \\ 0 & 1 \end{bmatrix},
        \sigma_{x} = X = \begin{bmatrix} 0 & 1 \\ 1 & 0 \end{bmatrix},
        \sigma_y = Y = \begin{bmatrix} 0 & -i \\ i & 0 \end{bmatrix},
        \sigma_z = Z = \begin{bmatrix} 1 & 0 \\ 0 & -1 \end{bmatrix}.
    \end{equation}
\end{remark}

We note the following useful expressions.
\begin{definition}[Bloch vectors]\label{rm:bloch_vector}
    \lean{Twirling.Physlib.blochState}\leanok
    For any single-qubit quantum state $\rho$, we can write it in terms of its Bloch vector,
    \begin{equation}
        \rho = \frac{1}{2} (I + r \cdot \sigma),
    \end{equation}
    where $r = (r_x, r_y, r_z) \in \mathbb{R}^3$ is the coordinate of the quantum state on the Bloch sphere and $\sigma = (\sigma_x, \sigma_y, \sigma_z)$ is a three-tuple of the Pauli operators.
\end{definition}

\begin{remark}[Pauli group identities]\label{rm:sigma_cubed}
    For all non-identity $\sigma_{i}, \sigma_{j}$ in the Pauli group where $i \neq j$, the following three statements are true.
    \begin{enumerate}
        \item $\sigma_{i}\sigma_{j} \sigma_{i} = - \sigma_j$,
        \item $ \sigma_{j}^3 = \sigma_{j}$, and
        \item $\sigma_i^2 = \sigma_i$, since $\sigma_i$ is represented by a self-adjoint matrix.
    \end{enumerate}
\end{remark}

We want to show the following simpler equivalence using matrix computations.
\begin{theorem}[Pauli twirling of a single qubit]\label{thm:twirl_single_qubit}
    \lean{Twirling.Physlib.twirl_single_qubit}\leanok
    \uses{rm:bloch_vector}
    For an arbitrary single-qubit state $\rho$,
    \begin{equation}
        \langle \rho \rangle_{G} = \frac{1}{N} \sum_{g} R(g) X R(g)^{\dagger} = \frac{I}{2},
    \end{equation}
    the maximally mixed state (since it is proportional to the identity matrix).
\end{theorem}

\begin{proof}\leanok
Note that $G$ has order $N = 4$.
For any Pauli operator $\pm (i) X \in G$, $X$ is self-adjoint.
Thus, we have that
\begin{equation}
    \pm (i) X \rho (\pm (i) X)^{\dagger} = (\pm i) (\mp i) X \rho X = X \rho X.
    \label{eq:simplified_sigma_rho_sigma*}
\end{equation}
By Def.~\ref{def:group_average} and Eq.~\ref{eq:simplified_sigma_rho_sigma*},
\begin{equation}
    \langle \rho \rangle_{G} = \frac{1}{4} \left( \sigma_{x} \rho \sigma_{x} + \sigma_{y} \rho \sigma_{y} + \sigma_{z} \rho \sigma_{z} + I \rho I \right).
\end{equation}
By Remark~\ref{rm:bloch_vector} and Remark~\ref{rm:sigma_cubed}, we can rewrite $\rho$ such that
\begin{equation}
    \begin{split}
        \langle \rho \rangle_{G} = \frac{1}{4} \left( \sigma_{x} (\frac{1}{2} (I + r \cdot \sigma) \sigma_{x} + \sigma_{y} (\frac{1}{2} (I + r \cdot \sigma)) \sigma_{y} + \sigma_{z} (\frac{1}{2} (I + r \cdot \sigma)) \sigma_{z} + I (\frac{1}{2} (I + r \cdot \sigma)) I \right)\\
        = \frac{1}{4} \left( \frac{1}{2} \left[ \sigma_x I \sigma_x + r_x \sigma_x^3 + r_y \sigma_x \sigma_y \sigma_x + r_z \sigma_x \sigma_z \sigma_x \right]
        + \frac{1}{2} \left[ \sigma_y I \sigma_y + r_x \sigma_y \sigma_x \sigma_y + r_y \sigma_y^3 + r_z \sigma_y \sigma_z \sigma_y \right]
        \right.\\
        + \frac{1}{2} \left[ \sigma_z I \sigma_z + r_x \sigma_z \sigma_x \sigma_z + r_y \sigma_z \sigma_y \sigma_z + r_z \sigma_z^3 \right]
        + \frac{1}{2} \left[ I^3 + r_x I \sigma_x I + r_y I \sigma_y I + r_z I \sigma_z I \right]
        \left. \right)\\
        = \frac{1}{8} \left( \right. [I + r_x \sigma_x - r_y \sigma_y - r_z \sigma_z]
        + [I - r_x \sigma_x + r_y \sigma_y - r_z \sigma_z]\\
        + [I - r_x \sigma_x - r_y \sigma_y + r_z \sigma_z]
        + [I + r_x \sigma_x + r_y \sigma_y + r_z \sigma_z ]\left. \right)\\
        = \frac{1}{8}\left( 4I\right) = \frac{I}{2}.
    \end{split}
\end{equation}
\end{proof}
